% !TeX root = ../thuthesis-example.tex

\chapter{自适应多步检索生成模型}

为实现真正“因材施教”的个性化方言教学，本项目设计了“自适应多步检索生成模型”。该模型的核心在于，它超越了传统检索增强生成（RAG）仅关注“信息完备性”的局限，创新性地引入了对“学习者理解度”的实时考量。
本项目提出的模型，是在RAG和CoRAG的坚实基础上，针对个性化教学场景进行的深度优化和创新。本项目的模型进行了两大根本性改造：
\begin{enumerate}
    \item \textbf{知识源的变革}：它并非从通用的网络文本中检索，而是以“\textbf{多模态方言音系知识图谱与现代汉语方言音档}” ($D_{\text{kg}}$) 作为其唯一的、权威的知识源泉。这使得检索过程从“大海捞针”转变为在结构化知识空间中的精准导航。
    \item \textbf{核心目标的转向}：模型的核心优化目标，从单纯的“信息完备性”转向了“\textbf{信息完备性}”与“\textbf{学习者理解度}”的动态平衡。它创新性地引入了对学习者状态的实时考量，从而构成一个信息-认知双重自适应的动态教学系统。
\end{enumerate}
本模型的整体概览如图 \ref{fig:corag_adapted} 所示。

\begin{figure}[H]
    \centering
    % 此处应放置您根据CoRAG图示改编的框架图
    \includegraphics[width=\textwidth]{coRAG.pdf}
    \caption{自适应多步检索生成概览图}
    \label{fig:corag_adapted}
\end{figure}

\section{学习者状态建模}
这是整个自适应机制的基石。本项目为每位用户 $u$ 构建一个动态更新的“学习者状态向量”$S_u$。该向量由两部分构成：$ S_u = \{ M_u, P_u \} $，其中，$M_u$ 是知识掌握度向量 (Mastery Vector)，$P_u$ 是认知偏好向量 (Preference Vector)。
\begin{itemize}
    \item \textbf{知识掌握度向量 $M_u$}: 该向量的维度与知识图谱中的核心知识节点一一对应，例如，$m_{u, \text{`入声韵'}} = 0.2$ 表示该用户对入声韵的掌握较差。
    \item \textbf{认知偏好向量 $P_u$}: 该向量量化了用户对不同信息呈现模态的偏好，例如，$p_{u, \text{`spectrogram'}} = 0.9$ 表明该学习者对视觉化的声谱图理解效率更高。
\end{itemize}
该模型 $S_u$ 在每次有效交互后都会通过一个更新函数 $S_{u, t+1} = F(S_{u, t}, \text{Interaction}_t)$ 进行实时更新。

\section{训练数据增强：生成自适应检索链}
为训练模型，本项目使用“\textbf{拒绝采样}” (Rejection Sampling) 思想来自动生成高质量的训练数据。过程如下：对于一个教学问题 $Q$ 和一个\textbf{模拟}的学习者状态 $S$，本项目生成多条候选的检索-生成链。与CoRAG仅评估答案似然不同，本项目设计了一个复合的\textbf{质量评估函数}来筛选最优教学链：
\begin{equation}
\text{Score}(\text{Chain}) = \alpha \log P(A | \text{Chain}, S) + \beta \text{Relevance}(\text{Chain}, S)
\end{equation}
其中，第一项代表“\textbf{信息完备性}”，第二项 $\text{Relevance}$ 代表“\textbf{个性化适应度}”。本项目选择得分最高的链作为一条高质量的训练实例 $(Q, S, A, Q_{1:L}, A_{1:L})$。

\section{多任务联合训练与复合损失函数}
在构建好增强型数据集后，本项目在一个多任务学习（multi-task learning）框架下，使用标准的下一词元预测（next-token prediction）目标对模型进行微调。此举旨在让模型同时掌握三种既独立又关联的核心能力：一是根据学习者状态和对话历史进行策略性子查询的能力；二是将检索到的知识图谱内容转化为适应个人理解能力的子答案的能力；三是整合整个交互过程的信息，生成最终个性化教学包的能力。

这三个任务由一个统一的复合损失函数 $\mathcal{L}_{\text{total}}$ 来共同优化，该函数是三个子任务损失的加权和。

\paragraph{子查询预测损失 ($\mathcal{L}_{\text{sub\_query}}$)}
子查询预测损失的目标是让模型学会根据当前对话历史和学习者状态 $S$ 的条件下，生成指向知识图谱的最合适的下一个子查询 $Q_i$。这个子查询可能是为了深入挖掘某个音系规则，也可能是为了寻找更简单的多媒体示例。其损失函数定义为标准的条件概率对数似然：
\begin{equation}
\mathcal{L}_{\text{sub\_query}} = - \log P(Q_i | Q, Q_{<i}, A_{<i}, S)
\end{equation}

其中，$Q$ 是初始问题，$Q_{<i}$ 和 $A_{<i}$ 是前 $i-1$ 步的子查询和子答案序列。关键在于，这个概率是以学习者状态 $S$ 为条件的。这意味着，对于同一个对话历史，面对一个初学者和一个进阶者，模型应该学会生成不同性质的子查询 $Q_i$。例如，模型可能会为初学者生成一个查询，用于在知识图谱中检索与某个错误发音相关的特定“音系规则”节点，或是查找“口型示范视频”等多模态资产。

\paragraph{子答案生成损失 ($\mathcal{L}_{\text{sub\_answer}}$)}
在生成每一步的子答案 $A_i$ 时，模型不仅要忠于从知识图谱 $D_{kg}$ 中检索到的内容 $D^{(i)}_{1:k}$，还必须根据状态 $S$ 调整其详略和风格。
\begin{equation}
\mathcal{L}_{\text{sub\_answer}} = - \log P(A_i | Q_i, D^{(i)}_{1:k}, S)
\end{equation}

在这里，$D^{(i)}_{1:k}$ 代表的“文档”是从方言音系知识图谱中检索出的结构化信息节点。这可能包括：一段描述某个变调规则的文本、一个包含标准音高曲线数据的节点、一个真人发音的音频文件，或是一个展示正确口型的视频片段。学习者状态 $S$ 的引入，使得模型在生成子答案 $A_i$ 时，能够自动调整其用词的复杂度、解释的详细程度以及举例的方式，使其内容既准确又易于理解。

\paragraph{最终答案生成损失 ($\mathcal{L}_{\text{final\_answer}}$)}
最终的答案 $A$ 是一个为用户量身定制的、丰富的“个性化教学包”。其生成过程受到最完整信息的约束，因此其损失函数定义为：
\begin{equation}
\mathcal{L}_{\text{final\_answer}} = - \log P(A | Q, Q_{1:L}, A_{1:L}, \mathcal{D}_{\text{all}}, S)
\end{equation}

其中，$\mathcal{D}_{\text{all}}$ 代表了在整个检索链中从知识图谱收集到的所有相关多模态资产与结构化知识。学习者状态 $S$ 在此的作用是决定最终答案的整体风格、结构和呈现方式（例如，是侧重于呈现图谱中的规则文本，还是突出音频和视频示范），从而将检索到的零散知识点，合成为一堂完整的、个性化的微型课程。

综上所述，模型的\textbf{总体训练目标}是最小化以下加权复合损失函数：
\begin{equation}
\mathcal{L}_{\text{total}} = w_{\text{query}} \cdot \mathcal{L}_{\text{sub\_query}} + w_{\text{answer}} \cdot \mathcal{L}_{\text{sub\_answer}} + w_{\text{final}} \cdot \mathcal{L}_{\text{final\_answer}}
\end{equation}

其中，$w_{\text{query}}$, $w_{\text{answer}}$, $w_{\text{final}}$ 是超参数，用于平衡三个子任务在训练过程中的重要性。通过优化这个总损失函数，本项目训练的不再是一个单纯的信息处理器，而是一个能够深度理解并策略性地利用多模态方言音系知识图谱，同时又能敏锐感知学习者状态、并动态调整自身教学策略的个性化教学智能体。

\section{推理时决策策略：最大化教学奖励与自适应教学路径生成}
在模型训练完成后，当面对一个带有真实状态向量 $S_u$ 的用户时，本项目需要一套高效的决策策略来实时生成最优的教学路径。这个决策过程可以被视为一个在复杂的“教学策略空间”中的启发式搜索问题，其目标是找到一条能最大化“\textbf{总教学奖励}”的路径。

本项目定义一个教学链（Chain）的总教学奖励 $\mathcal{R}_{\text{pedagogical}}$ 为以下三项的加权和：
\begin{equation}
\mathcal{R}_{\text{pedagogical}}(\text{Chain}) = w_{\text{correct}} \cdot \mathcal{R}_{\text{correct}} + w_{\text{adapt}} \cdot \mathcal{R}_{\text{adapt}} + w_{\text{effic}} \cdot \mathcal{R}_{\text{effic}}
\end{equation}
其中，$w$为权重超参数。

\begin{enumerate}
    \item \textbf{正确性奖励 ($\mathcal{R}_{\text{correct}}$)}: 这是最基础的奖励，衡量该教学链最终能否引导模型生成正确、符合事实的答案 $A$。本项目用模型在给定该链后，生成标准答案的对数似然来近似它：
    \begin{equation}
        \mathcal{R}_{\text{correct}} = \log P(A | \text{Chain}, S_u)
    \end{equation}
    
    \item \textbf{适应性奖励 ($\mathcal{R}_{\text{adapt}}$)}: 这是实现“因材施教”的核心。该奖励衡量教学链中的内容（知识的难度、呈现的模态）与当前学习者状态 $S_u = \{M_u, P_u\}$ 的匹配度。例如，它可以被定义为链中所有步骤的适应性得分之和：
    \begin{equation}
        \mathcal{R}_{\text{adapt}} = \sum_{i \in \text{Chain}} \left( (1 - \text{Complexity}(A_i)) \cdot (1 - m_{u,j}) + \text{ModalityMatch}(A_i, P_u) \right)
    \end{equation}
    其中，$\text{Complexity}(A_i)$是子答案$A_i$的复杂度，当它与用户的掌握度$m_{u,j}$都较低时，奖励较高。$\text{ModalityMatch}$则衡量子答案$A_i$使用的模态是否符合用户的偏好$P_u$。
    
    \item \textbf{效率奖励 ($\mathcal{R}_{\text{effic}}$)}: 该奖励惩罚冗长、低效的教学路径，确保教学过程的简洁性。它可以简单地用教学链长度 $L$ 的倒数来表示：
    \begin{equation}
        \mathcal{R}_{\text{effic}} = \frac{1}{L}
    \end{equation}
\end{enumerate}

基于上述教学奖励函数，本项目实现了CoRAG中提出的多种解码策略的自适应版本：

\begin{itemize}
    \item \textbf{贪心解码 (Greedy Decoding)}: 在每一步都选择能使\textbf{单步}预期奖励最大化的子查询。此策略速度最快，但忽略了长远收益。
    \item \textbf{Best-of-N 采样 (Best-of-N Sampling)}: 通过采样生成 $N$ 条完整的候选教学链，然后使用上述定义的 $\mathcal{R}_{\text{pedagogical}}$ 对每条链进行\textbf{完整评估}，选择总奖励最高的链条。
    \item \textbf{树搜索 (Tree Search)}: 在构建决策树时，每个节点的扩展都基于对未来路径的“\textbf{预期总教学奖励}”的估算。这使得模型能够进行更具前瞻性的、全局最优的教学规划。
\end{itemize}
\section{多模态整合与闭环教学}
在最优的教学链被选定后，系统便进入了将检索结果转化为最终教学体验的关键阶段。此阶段的核心任务是整合该链条上从知识图谱中检索出的所有信息资产，生成一个“个性化教学包”（Personalized Teaching Package）。更重要的是，这个教学包不仅是本次交互的终点，更是下一轮自适应学习的起点，从而形成一个“\textbf{诊断-教学-评估-再诊断}”的完整闭环。

\begin{figure}[H]
    \centering
    % 此处应放置您根据CoRAG图示改编的框架图
    \includegraphics[scale=0.4]{个性化教学.pdf}
    \caption{多模态整合与闭环教学流程图}
    \label{fig:多模态整合与闭环教学}
\end{figure}

本项目定义的“个性化教学包”并非一段简单的文本，而是一个结构化的、待渲染的数据对象（例如一个JSON对象）。该对象由推理聚合智能体（Manager Agent）依据最优教学链生成，其内容主要包含以下几个模块：

\paragraph{诊断摘要 (Diagnostic Summary)} 
由大语言模型生成的、针对本次用户输入（如一次错误发音）的、高度定制化的文本解释。例如：“您的‘入’字发音中，主要问题在于韵母的口腔闭合过快，导致其听起来像一个短促的元音，而忽略了它作为入声韵尾所特有的喉塞音[-p̚]、[-t̚]或[-k̚]。”

\paragraph{核心教学模块 (Core Remediation Nodes)}
这是一个包含了多个教学“卡片”的列表。每一张卡片都围绕一个从知识图谱中检索出的核心知识点，并根据学习者状态$S_u$来决定其呈现方式。常见的卡片类型有：
\begin{enumerate}
    \item \textbf{规则讲解卡}: 以文本形式呈现从知识图谱中提取的音系规则，但其用词和详略程度会根据用户掌握度$M_u$动态调整。
    \item \textbf{多模态示范卡}: 这是本框架的核心。它将多种模态的资产进行并列或叠加展示，以提供最直观的感知体验。例如，一张“音高曲线对比卡”可以同时展示标准发音和用户发音的音高曲线图，并允许用户点击播放两条曲线对应的音频，形成视觉和听觉上的直接对比。一张“口型示范卡”则可能包含标准口型的前视图、侧视图慢放视频。
    \item \textbf{互动练习卡}: 提供新的、与当前教学知识点相关的练习题。例如，在讲解了“连续变调”后，提供两个新的双字词让用户进行跟读练习。
\end{enumerate}

\paragraph{知识图谱链接 (KG Links)}
在教学包的关键术语（如“入声”、“喉塞音”）旁，会嵌入指向知识图谱中对应节点的超链接，鼓励有余力的用户进行更深层次的探索式学习。

值得注意的是，大语言模型负责生成这个结构化的数据对象（即内容逻辑），而最终用户看到的、包含音频播放器、可交互图表的精美界面，则是由一个独立的“\textbf{前端渲染引擎}” (Frontend Rendering Engine) 负责解析该数据对象并呈现的。这种逻辑与呈现分离的架构，使得教学内容的迭代和展示形式的优化可以独立进行。

\subsubsection*{闭环反馈：从交互到状态演化}
“个性化教学包”不仅是知识的输出，更是新一轮数据采集的“传感器阵列”。用户的每一个微观交互行为都会被前端的事件监听器捕捉，并作为事件$I_t$发送回后端。

\paragraph{交互事件捕捉}
捕捉的事件类型包括但不限于：`play\_audio(asset\_uri)`、`replay\_video(asset\_uri)`、`zoom\_in\_spectrogram(asset\_uri)`、`click\_kg\_link(node\_uri)`，以及最重要的`submit\_exercise(user\_audio\_uri, target\_word)`。

\paragraph{学习者状态演化}
自适应反馈智能体根据接收到的事件$I_t$，通过一个更新函数 $S_{u, t+1} = F(S_{u, t}, I_t)$ 来更新学习者状态向量。

\textbf{掌握度向量$M_u$的更新}: 当接收到练习提交事件时，系统会调用一个发音评估模块（如基于ASR的评分模型）对用户的发音进行打分（$\text{score} \in [0,1]$）。该分数将通过一个带有学习率$\eta$的阻尼更新规则，来更新$M_u$中对应知识节点的掌握度：
$$ m_{u,j}^{\text{new}} = m_{u,j}^{\text{old}} + \eta \cdot (\text{score} - m_{u,j}^{\text{old}}) $$
这使得用户的知识画像能够根据其实际表现进行平滑的演进。

\textbf{认知偏好向量$P_u$的更新}: 当用户频繁地与某一类型的多媒体资产交互时（例如，反复播放音频，或长时间查看谱图），系统会将其解读为用户对该模态的偏好或依赖。相应地，该模态在认知偏好向量$P_u$中的权重会获得少量提升，随后整个向量被重新归一化。

通过这一机制，更新后的状态向量$S_{u, t+1}$便成为了用户下一次提问或练习时的“初始状态”。系统对用户的理解因此变得更加精准，从而能够在下一轮交互中，生成一个更加贴合其当前水平和偏好的、全新的个性化教学包，最终形成一个不断螺旋上升的自适应教学闭环。


